% Generated draft: requires mathematical and source review.
\documentclass[11pt]{article}
\usepackage[margin=25mm]{geometry}
\usepackage{fontspec}
\setmainfont{Latin Modern Roman}
\usepackage{amsmath,amssymb,mathtools,unicode-math}
\setmathfont{Latin Modern Math}
\usepackage{enumitem,needspace,xurl,hyperref}
\setlength{\parindent}{0pt}
\setlength{\parskip}{5pt}
\setlength{\emergencystretch}{3em}
\newcommand{\N}{\mathbb{N}}
\newcommand{\Z}{\mathbb{Z}}
\newcommand{\Q}{\mathbb{Q}}
\newcommand{\R}{\mathbb{R}}
\newcommand{\C}{\mathbb{C}}
\newcommand{\F}{\mathbb{F}}
\newcommand{\A}{\mathbb{A}}
\newcommand{\PP}{\mathbb{P}}
\newcommand{\E}{\mathbb{E}}
\newcommand{\eps}{\varepsilon}
\newcommand{\dd}{\,\mathrm{d}}
\begin{document}
\section*{UnsolvedMath ID: 30004538}
\textbf{Source problem code:} catalogue entry
\subsection{Definitions and mathematical statement}
The supplied record describes a research programme rather than a single quantified conjecture. The following notation makes its questions explicit without supplying missing model hypotheses.

Under the real-valued, continuous-time interpretation suggested by persistence on an interval, let \(X=(X_t)_{t\in\R}\) be a stationary Gaussian process: every finite vector of its coordinates is jointly Gaussian, and its finite-dimensional distributions are invariant under time translation. No centering, variance normalization, covariance function, spectral measure, or sample-path regularity is specified in the source.

For \(\ell\in\R\) and \(T>0\), distinguish the events
\[
 A^{>}_{T,\ell}=\{X_t>\ell\text{ for every }t\in[0,T]\},
 \qquad
 A^{\ge}_{T,\ell}=\{X_t\ge\ell\text{ for every }t\in[0,T]\}.
\]
For a model and process version in which the chosen event is measurable, write
\[
 p_X^{\diamond}(T,\ell)=\mathbb P(A^{\diamond}_{T,\ell}),
 \qquad \diamond\in\{>,\ge\}.
\]
A conventional exponential persistence exponent is
\[
 \theta_X^{\diamond}(\ell)
 =-\lim_{T\to\infty}\frac{1}{T}\log p_X^{\diamond}(T,\ell),
\]
provided the limit exists, with \(-\log 0=+\infty\). The source does not specify this normalization, guarantee existence of the limit, or settle whether strict or non-strict persistence is intended; the two boundary conventions must therefore remain distinct.

Whenever \(p_X^{\diamond}(T,\ell)>0\), the conditioned process is described by
\[
 \mathcal L\bigl((X_t)_{0\le t\le T}\mid A^{\diamond}_{T,\ell}\bigr).
\]
The supplied questions are to compute persistence exponents in specific stationary Gaussian examples, determine the typical shape under this conditional law, and investigate extensions to non-Gaussian or non-stationary processes. No particular examples or extension classes are identified. An asymptotic interpretation of ``typical shape'' would require a specified limiting regime, time and height scalings, path space, and mode of convergence or concentration; none is supplied. Thus the record does not determine a unique limiting-shape assertion or a universal extension theorem.

The supplied dated curation labels the programme partially solved and reports results for conditioned shapes and extensions in several spectral-singularity classes, while describing the general shape and extension questions as open. The relevant classes and hypotheses are unspecified, and these literature claims are not independently verified here.
\subsection{Short English statement}
For stationary Gaussian processes, compute persistence exponents in concrete examples and describe paths conditioned to stay above a level \(\ell\) throughout \([0,T]\). Investigate analogous questions for non-Gaussian or non-stationary models. The source leaves the examples, boundary convention, exponent normalization, and meaning of typical shape unspecified. Its partially solved designation is supplied dated curation.
\subsection{Sources}
\begin{itemize}
\item Supplied local UnsolvedMath registry record, permanent ID \(30004538\), title \emph{Values and Conditioned Shapes of Gaussian Persistence}, source problem code \texttt{catalogue entry}, in the Oberwolfach Reports Open Problems collection. No individual proposer is supplied.
\item Supplied report citation: \emph{Stochastic Processes under Constraints} (2021), Oberwolfach Reports, DOI 10.4171/owr/2020/32, pp. 1601--1656. \url{https://doi.org/10.4171/owr/2020/32}. The background describes Oberwolfach Reports as documenting workshops and their open-problem discussions. The record separately gives proposed year 2020; the relation between these year labels is not resolved here.
\item Supplied literature reference: \emph{Gaussian persistence exponents}. \url{https://arxiv.org/abs/2112.04820}.
\item Supplied literature reference: \emph{Conditioned field asymptotics}. \url{https://arxiv.org/abs/2605.20587}.
\item Literature assessment dated exactly \texttt{2026-08-22}: supplied dated curation, not independently verified. It labels the problem partially solved, reports conditioned persistence shapes and extensions beyond the basic stationary Gaussian setting for several spectral-singularity classes, and states that typical shapes and non-Gaussian or non-stationary generalizations in full generality remain open. The supplied background requests independent verification of time-sensitive status and cited claims.
\item The registry's \texttt{corrected\_verified} statement label is accompanied by an explanation that only formula delimiters, notation, grammar, and typography were repaired, with mathematical content and scope preserved. It is not treated here as independent mathematical or bibliographic verification.
\end{itemize}
\end{document}
